\section{THỐNG KÊ MẪU SỐ LIỆU GHÉP NHÓM}
\subsection{\vdmh}
\Opensolutionfile{ans}[ans/DA-CD5-THONGKE-VD]
\begin{vd}
    Trong các số đặc trưng đo xu thế trung tâm dưới đây, số nào thỏa mãn có $25\%$ giá trị trong mẫu số liệu nhỏ hơn nó và $75\%$ giá trị trong mẫu số liệu lớn hơn nó?
    \choice
    {Số trung bình}
    {Trung vị}
    {\True Tứ phân vị thứ nhất}
    {Tứ phân vị thứ ba}
    %    \loigiai{
        %        Tứ phân vị thứ nhất ($Q_1$) là giá trị chia mẫu số liệu (đã sắp xếp) sao cho có khoảng $25\%$ số liệu nhỏ hơn và $75\%$ số liệu lớn hơn nó.
        %    }
\end{vd}
\begin{vd}[MH BGD-2025]
    Hai mẫu số liệu ghép nhóm $M_1, M_2$ có bảng tần số ghép nhóm như sau:
    \begin{center}
        $M_1$ \quad
    \begin{tabular}{|c|c|c|c|c|c|}
        \hline
        Nhóm & $[8;10)$ & $[10;12)$ & $[12;14)$ & $[14;16)$ & $[16;18)$ \\
        \hline
        Tần số & $3$ & $4$ & $8$ & $6$ & $4$ \\
        \hline
    \end{tabular}
    \\
    $M_2$ \quad
    \begin{tabular}{|c|c|c|c|c|c|}
        \hline
        Nhóm & $[8;10)$ & $[10;12)$ & $[12;14)$ & $[14;16)$ & $[16;18)$ \\
        \hline
        Tần số & $6$ & $8$ & $16$ & $12$ & $8$ \\
        \hline
    \end{tabular}
    \end{center}
    Gọi $s_1, s_2$ lần lượt là độ lệch chuẩn của mẫu số liệu ghép nhóm $M_1, M_2$. Phát biểu nào sau đây là đúng?
    \choice
    {\True $s_1 = s_2$}
    {$s_1 = 2s_2$}
    {$2s_1 = s_2$}
    {$4s_1 = s_2$}
%    \loigiai{
%        Ta nhận thấy tần số của mỗi nhóm trong mẫu $M_2$ đều gấp đôi tần số của nhóm tương ứng trong mẫu $M_1$. Do đó, tỉ lệ phần trăm (tần suất) của mỗi nhóm tương ứng trong hai mẫu số liệu là như nhau.
%        Vì tần suất phân bố như nhau nên số trung bình và phương sai (kéo theo độ lệch chuẩn) của hai mẫu số liệu này là bằng nhau.
%        Suy ra $s_1 = s_2$.
%    }
\end{vd}
\begin{vd}[BGD-2025]
    Một người chia thời lượng (đơn vị: giây) thực hiện các cuộc gọi điện thoại của mình trong một tuần thành sáu nhóm và lập bảng tần số ghép nhóm như sau:
    \begin{center}
        \begin{tabular}{|c|c|c|c|c|c|c|}
            \hline
            Nhóm & $[0;30)$ & $[30;60)$ & $[60;90)$ & $[90;120)$ & $[120;150)$ & $[150;180)$ \\
            \hline
            Tần số & $10$ & $11$ & $8$ & $6$ & $4$ & $1$ \\
            \hline
        \end{tabular}
    \end{center}
    Tứ phân vị thứ ba $Q_3$ (đơn vị: giây) của mẫu số liệu ghép nhóm trên bằng
    \choice
    {$100$}
    {\True $95$}
    {$105$}
    {$90$}
    \loigiai{
        Cỡ mẫu $n = 10 + 11 + 8 + 6 + 4 + 1 = 40$. \\
        Vị trí của tứ phân vị thứ ba là $\frac{3n}{4} = \frac{3 \cdot 40}{4} = 30$. \\
        Tần số tích lũy của các nhóm lần lượt là: $10; 21; 29; 35; 39; 40$. \\
        Do đó, nhóm chứa tứ phân vị thứ ba là $[90;120)$. \\
        Áp dụng công thức tính tứ phân vị, ta có:
        $$Q_3 = 90 + \frac{30 - 29}{6} \cdot (120 - 90) = 90 + \frac{1}{6} \cdot 30 = 95.$$
    }
\end{vd}
\begin{vd}[Chuyên Lương Thế Vinh-Đồng Nai-2025] Mỗi ngày bác Hương đều đi bộ để rèn luyện sức khỏe. Quãng đường đi bộ mỗi ngày (đơn vị km) của bác Hương trong 20 ngày được thống kê lại ở bảng sau:\\
    \centerline{
        \begin{tabular}{|c|c|c|c|c|c|}
            \hline
            Quãng đường (km) & $[2,7;3,0)$ & $[3,0;3,3)$ & $[3,3;3,6)$ & $[3,6;3,9)$ & $[3,9;4,2)$ \\
            \hline
            Số ngày & $3$ & $6$ & $5$ & $4$ & $2$ \\
            \hline
        \end{tabular}
    }
    Khoảng tứ phân vị của mẫu số liệu ghép nhóm là
    \choice
    {$0,9$}
    {$0,975$}
    {$0,5$}
    {\True $0,575$}
    \loigiai{
        Cỡ mẫu $n = 20$.
        \begin{itemize}
            \item Vị trí tứ phân vị thứ nhất là $\frac{n}{4} = \frac{20}{4} = 5$. Tần số tích lũy các nhóm đầu: $3; 9$. Suy ra nhóm chứa $Q_1$ là $[3,0;3,3)$.\\
            $$Q_1 = 3,0 + \frac{5 - 3}{6} \cdot (3,3 - 3,0) = 3,0 + \frac{2}{6} \cdot 0,3 = 3,1.$$
            \item Vị trí tứ phân vị thứ ba là $\frac{3n}{4} = \frac{3 \cdot 20}{4} = 15$. Tần số tích lũy các nhóm đầu: $3; 9; 14; 18$. Suy ra nhóm chứa $Q_3$ là $[3,6;3,9)$.\\
            $$Q_3 = 3,6 + \frac{15 - 14}{4} \cdot (3,9 - 3,6) = 3,6 + \frac{1}{4} \cdot 0,3 = 3,675.$$
        \end{itemize}
        Khoảng tứ phân vị của mẫu số liệu ghép nhóm là: $\Delta_Q = Q_3 - Q_1 = 3,675 - 3,1 = 0,575$.
    }
    \dotlineEX{2}
\end{vd}
\begin{vd}
    Bạn Chi rất thích nhảy hiện đại. Thời gian tập nhảy mỗi ngày trong thời gian gần đây của bạn Chi được thống kê lại ở bảng sau:
    \begin{center}
        \begin{tabular}{|l|c|c|c|c|c|}
            \hline
            Thời gian (phút) & $[20; 25)$ & $[25; 30)$ & $[30; 35)$ & $[35; 40)$ & $[40; 45)$ \\ \hline
            Số ngày & $6$ & $6$ & $4$ & $1$ & $1$ \\ \hline
        \end{tabular}
    \end{center}
    \choiceTF
    {\True Giá trị đại diện của nhóm $[30; 35)$ là $32,5$}
    {\True Mốt của mẫu số liệu ghép nhóm này bằng $25$}
    {Trung vị $Q_2$ của mẫu số liệu này bằng $28$}
    {\True Có $75\%$ số ngày bạn Chi tập luyện không quá $31,875$ phút}
    \loigiai{
        Cỡ mẫu $n = 6 + 6 + 4 + 1 + 1 = 18$.
        \begin{itemize}
            \item \textbf{Ý a):} Giá trị đại diện của nhóm $[30; 35)$ là $x_3 = \frac{30+35}{2} = 32,5$. Do đó ý a) \textbf{Đúng}.
            \item \textbf{Ý b):} Tần số lớn nhất là $6$ tại hai nhóm $[20; 25)$ và $[25; 30)$. 
            \begin{itemize}
                \item Nếu tính mốt cho nhóm $[20; 25)$: $M_o = 20 + \frac{6-0}{(6-0) + (6-6)} \cdot 5 = 25$.
                \item Nếu tính mốt cho nhóm $[25; 30)$: $M_o = 25 + \frac{6-6}{(6-6) + (6-4)} \cdot 5 = 25$.
            \end{itemize}
            Vậy mốt của mẫu số liệu là $25$. Do đó ý b) \textbf{Đúng}.
            \item \textbf{Ý c):} Nhóm chứa trung vị là $[25; 30)$. 
            Ta có: $Q_2 = 25 + \frac{\frac{18}{2} - 6}{6} \cdot 5 = 27,5$. Do đó ý c) \textbf{Sai}.
            \item \textbf{Ý d):} Tứ phân vị thứ ba $Q_3$ thuộc nhóm $[30; 35)$.
            Ta có: $Q_3 = 30 + \frac{13,5 - 12}{4} \cdot 5 = 31,875$.
            Ý nghĩa: Khoảng $75\%$ giá trị trong mẫu số liệu nhỏ hơn hoặc bằng $31,875$. Do đó ý d) \textbf{Đúng}.
        \end{itemize}
    }
\end{vd}
\begin{vd}%[Câu 8][857086][MĐ2]
    Cho mẫu số liệu dưới dạng bảng sau
    \begin{center}
        \begin{tabular}{|c|c|c|c|c|c|c|}
            \hline
            Số câu trả lời đúng & $[16;21)$ & $[21;26)$ & $[26;31)$ & $[31;36)$ & $[36;41)$ & \\
            \hline
            Tần số & $4$ & $6$ & $8$ & $18$ & $4$ & $N = 40$ \\
            \hline
        \end{tabular}
    \end{center}
    \choiceTF
    {\True Giá trị đại diện của lớp $[36;41)$ là $38,5$}
    {\True Công thức tính số trung bình là $\overline{x} = \dfrac{18,5.4 + 23,5.6 + 28,5.8 + 33,5.18 + 38,5.4}{40}$}
    {\True Số trung bình là $30$}
    {\True Phương sai của mẫu số liệu là $S^2 = 32,75$}
    \loigiai{
        \begin{itemize}
            \item a) Đúng. $\frac{36+41}{2} = 38,5$.
            \item b) Đúng. Công thức áp dụng chuẩn xác giá trị đại diện và tần số.
            \item c) Đúng. Bấm máy tính theo công thức ở ý b thu được $30$.
            \item d) Đúng. $s^2 = \frac{18,5^2 \cdot 4 + \dots + 38,5^2 \cdot 4}{40} - 30^2 = 32,75$.
        \end{itemize}
    }
\end{vd}
\begin{vd}
    Thời gian hoàn thành một bài viết chính tả của một số học sinh lớp 4 tại hai trường $X$ và $Y$ được ghi lại ở bảng sau:
    \begin{center}
        \begin{tabular}{|l|c|c|c|c|c|}
            \hline
            Thời gian (phút) & $[6; 7)$ & $[7; 8)$ & $[8; 9)$ & $[9; 10)$ & $[10; 11)$ \\ \hline
            Số học sinh trường $X$ & $8$ & $10$ & $13$ & $10$ & $9$ \\ \hline
            Số học sinh trường $Y$ & $4$ & $12$ & $17$ & $14$ & $3$ \\ \hline
        \end{tabular}
    \end{center}
    \choiceTF
    {\True Thời gian hoàn thành bài viết trung bình của học sinh trường $X$ lâu hơn so với học sinh trường $Y$}
    {\True Khoảng biến thiên của mẫu số liệu ghép nhóm về thời gian hoàn thành bài viết của học sinh hai trường là bằng nhau}
    {Khoảng tứ phân vị của mẫu số liệu trường $Y$ xấp xỉ bằng $2,2$}
    {\True Học sinh trường $Y$ có thời gian hoàn thành bài viết đồng đều hơn học sinh trường $X$}
    \loigiai{
        Cỡ mẫu của cả hai trường là $n_X = n_Y = 50$. 
        Giá trị đại diện các nhóm lần lượt là: $x_1=6,5; x_2=7,5; x_3=8,5; x_4=9,5; x_5=10,5$.
        \begin{itemize}
            \item \textbf{Ý a):} 
            Trung bình trường $X$: $\bar{x}_X = \frac{8 \cdot 6,5 + 10 \cdot 7,5 + 13 \cdot 8,5 + 10 \cdot 9,5 + 9 \cdot 10,5}{50} = 8,54$. \\
            Trung bình trường $Y$: $\bar{x}_Y = \frac{4 \cdot 6,5 + 12 \cdot 7,5 + 17 \cdot 8,5 + 14 \cdot 9,5 + 3 \cdot 10,5}{50} = 8,5$. \\
            Vì $8,54 > 8,5$ nên ý a) \textbf{Đúng}.
            \item \textbf{Ý b):} Khoảng biến thiên của mẫu số liệu ghép nhóm được tính bằng hiệu giữa đầu mút phải của nhóm cuối cùng và đầu mút trái của nhóm đầu tiên. 
            $R_X = R_Y = 11 - 6 = 5$. Do đó ý b) \textbf{Đúng}.
            \item \textbf{Ý c):} Đối với trường $Y$:
            $Q_1$ thuộc nhóm $[7; 8)$ (vì $4 < 12,5 < 16$): $Q_1 = 7 + \frac{12,5 - 4}{12} \cdot 1 \approx 7,71$. \\
            $Q_3$ thuộc nhóm $[9; 10)$ (vì $33 < 37,5 < 47$): $Q_3 = 9 + \frac{37,5 - 33}{14} \cdot 1 \approx 9,32$. \\
            Khoảng tứ phân vị trường $Y$: $\Delta_{Q_Y} = Q_3 - Q_1 \approx 1,61$. (Giá trị $2,2$ là của trường $X$). Do đó ý c) \textbf{Sai}.
            \item \textbf{Ý d):} Tính phương sai và độ lệch chuẩn:
            $s_X^2 = \frac{1}{50} \sum n_i x_i^2 - (\bar{x}_X)^2 \approx 1,76 \Rightarrow s_X \approx 1,33$. \\
            $s_Y^2 = \frac{1}{50} \sum n_i x_i^2 - (\bar{x}_Y)^2 = 1,08 \Rightarrow s_Y \approx 1,04$. \\
            Vì $s_Y < s_X$ nên độ phân tán của trường $Y$ thấp hơn, tức là học sinh trường $Y$ hoàn thành bài viết đồng đều hơn. Do đó ý d) \textbf{Đúng}.
        \end{itemize}
    }
\end{vd}
\Closesolutionfile{ans}
\subsection{\btrl}
\Opensolutionfile{ans}[ans/DA-CD5-THONGKE-BTRL]
\begin{ex}
    Trong các số đặc trưng đo xu thế trung tâm dưới đây, số nào thỏa mãn có $25\%$ giá trị trong mẫu số liệu nhỏ hơn nó và $75\%$ giá trị trong mẫu số liệu lớn hơn nó?
    \choice
    {Số trung bình}
    {Trung vị}
    {\True Tứ phân vị thứ nhất}
    {Tứ phân vị thứ ba}
%    \loigiai{
%        Tứ phân vị thứ nhất ($Q_1$) là giá trị chia mẫu số liệu (đã sắp xếp) sao cho có khoảng $25\%$ số liệu nhỏ hơn và $75\%$ số liệu lớn hơn nó.
%    }
\end{ex}
\begin{ex}
    Trong các số đặc trưng đo xu thế trung tâm dưới đây, số nào thỏa mãn có $75\%$ giá trị trong mẫu số liệu nhỏ hơn nó và $25\%$ giá trị trong mẫu số liệu lớn hơn nó?
    \choice
    {Số trung bình}
    {Trung vị}
    {Tứ phân vị thứ nhất}
    {\True Tứ phân vị thứ ba}
%    \loigiai{
%        Tứ phân vị thứ ba ($Q_3$) là giá trị chia mẫu số liệu (đã sắp xếp) sao cho có khoảng $75\%$ số liệu nhỏ hơn và $25\%$ số liệu lớn hơn nó.
%    }
\end{ex}
\begin{ex}
    Khảo sát thời gian tập thể dục của một số học sinh khối 11 thu được mẫu số liệu ghép nhóm sau:\\
    \centerline{\begin{tabular}{|c|c|c|c|c|c|}
            \hline
            Thời gian (phút) & $[0;20)$ & $[20;40)$ & $[40;60)$ & $[60;80)$ & $[80;100)$ \\
            \hline
            Số học sinh & $5$ & $9$ & $12$ & $10$ & $6$ \\
            \hline
        \end{tabular}}
    Mẫu số liệu ghép nhóm đã cho có tất cả bao nhiêu nhóm?
    \choice
    {$3$}
    {$4$}
    {\True $5$}
    {$6$}
%    \loigiai{
%        Nhìn vào bảng, ta đếm được có $5$ nhóm số liệu.
%    }
\end{ex}
\begin{ex}
    Khảo sát thời gian tập thể dục của một số học sinh khối 11 thu được mẫu số liệu ghép nhóm sau:\\
    \centerline{
        \begin{tabular}{|c|c|c|c|c|c|}
            \hline
            Thời gian (phút) & $[0;20)$ & $[20;40)$ & $[40;60)$ & $[60;80)$ & $[80;100)$ \\
            \hline
            Số học sinh & $5$ & $9$ & $12$ & $10$ & $6$ \\
            \hline
        \end{tabular}
    }
    Nhóm chứa mốt của mẫu số liệu ghép nhóm trên là
    \choice
    {$[20;40)$}
    {\True $[40;60)$}
    {$[60;80)$}
    {$[80;100)$}
%    \loigiai{
%        Tần số lớn nhất là $12$, tương ứng với nhóm $[40;60)$. Vậy nhóm chứa mốt là $[40;60)$.
%    }
\end{ex}
\begin{ex}
    Mức thưởng tết (triệu đồng) cho các nhân viên của một công ty được thống kê trong bảng sau:\\
    \centerline{
        \begin{tabular}{|c|c|c|c|c|c|}
            \hline
            Mức thưởng tết & $[5;10)$ & $[10;15)$ & $[15;20)$ & $[20;25)$ & $[25;30)$ \\
            \hline
            Số nhân viên & $13$ & $35$ & $47$ & $25$ & $10$ \\
            \hline
        \end{tabular}
    }
    Có bao nhiêu nhân viên trong công ty nhận được mức thưởng tết từ $15$ triệu đồng đến dưới $20$ triệu đồng?
    \choice
    {$5$}
    {$13$}
    {\True $47$}
    {$130$}
%    \loigiai{
%        Nhìn vào bảng, nhóm $[15;20)$ có tần số là $47$.
%    }
\end{ex}
\begin{ex}
    Mức thưởng tết (triệu đồng) cho các nhân viên của một công ty được thống kê trong bảng sau:\\
    \centerline{
        \begin{tabular}{|c|c|c|c|c|c|}
            \hline
            Mức thưởng tết & $[5;10)$ & $[10;15)$ & $[15;20)$ & $[20;25)$ & $[25;30)$ \\
            \hline
            Số nhân viên & $13$ & $35$ & $47$ & $25$ & $10$ \\
            \hline
        \end{tabular}
    }
    Số mốt của mẫu số liệu ghép nhóm trên là
    \choice
    {$0$}
    {\True $1$}
    {$2$}
    {$3$}
%    \loigiai{
%        Bảng tần số chỉ có một nhóm có tần số lớn nhất là nhóm $[15;20)$ (với tần số 47). Do đó, mẫu số liệu có $1$ mốt.
%    }
\end{ex}
\begin{ex}
    Khảo sát về số giờ mượn sách thư viện của học sinh khối 11 trường Y ta được mẫu số liệu ghép nhóm như sau:\\
    \centerline{
        \begin{tabular}{|c|c|c|c|c|c|}
            \hline
            Số giờ mượn & $[0;4)$ & $[4;8)$ & $[8;12)$ & $[12;16)$ & $[16;20)$ \\
            \hline
            Số học sinh & $54$ & $78$ & $120$ & $45$ & $12$ \\
            \hline
        \end{tabular}
    }
    Mốt của mẫu số liệu ghép nhóm trên gần nhất với giá trị nào sau đây?
    \choice
    {$120$}
    {$12$}
    {$8$}
    {\True $9$}
%    \loigiai{
%        Nhóm chứa mốt là $[8;12)$ vì có tần số lớn nhất $n_3 = 120$.
%        Áp dụng công thức tính mốt: 
%        $M_o = 8 + \frac{120 - 78}{(120 - 78) + (120 - 45)} \cdot (12 - 8) = 8 + \frac{42}{42 + 75} \cdot 4 \approx 9,43$.
%        Giá trị gần nhất là $9$.
%    }
\dotlineEX{2}
\end{ex}
\begin{ex}
    Khảo sát về số giờ mượn sách thư viện của học sinh khối 11 trường Y ta được mẫu số liệu ghép nhóm như sau:
    \begin{center}
        \begin{tabular}{|c|c|c|c|c|c|}
            \hline
            Số giờ mượn & $[0;4)$ & $[4;8)$ & $[8;12)$ & $[12;16)$ & $[16;20)$ \\
            \hline
            Số học sinh & $54$ & $78$ & $120$ & $45$ & $12$ \\
            \hline
        \end{tabular}
    \end{center}
    Số trung vị của mẫu số liệu trên là
    \choice
    {\True $8,75$}
    {$8$}
    {$10$}
    {$9$}
    \loigiai{
        Cỡ mẫu $N = 54 + 78 + 120 + 45 + 12 = 309$.
        Ta có $\frac{N}{2} = 154,5$. Tần số tích lũy của nhóm 2 là $54 + 78 = 132$, của nhóm 3 là $132 + 120 = 252$.
        Do đó trung vị thuộc nhóm $[8;12)$.
        $M_e = 8 + \frac{154,5 - 132}{120} \cdot 4 = 8,75$.
    }
\end{ex}
\begin{ex}
    Khảo sát về cân nặng của các học sinh lớp 11D3 người ta được mẫu dữ liệu ghép nhóm như sau:
    \begin{center}
        \begin{tabular}{|c|c|c|c|c|c|c|}
            \hline
            Cân nặng & $[30;40)$ & $[40;50)$ & $[50;60)$ & $[60;70)$ & $[70;80)$ & $[80;90)$ \\
            \hline
            Số học sinh & $2$ & $10$ & $16$ & $8$ & $2$ & $2$ \\
            \hline
        \end{tabular}
    \end{center}
    Tứ phân vị thứ nhất của mẫu số liệu trên là
    \choice
    {$50$}
    {\True $48$}
    {$40$}
    {$45$}
    \loigiai{
        Cỡ mẫu $N = 2 + 10 + 16 + 8 + 2 + 2 = 40$. Vị trí của $Q_1$ là $\frac{N}{4} = 10$.
        Tần số tích lũy của nhóm 1 là $2$, của nhóm 2 là $12$. Nhóm chứa $Q_1$ là $[40;50)$.
        $Q_1 = 40 + \frac{10 - 2}{10} \cdot 10 = 48$.
    }
\end{ex}
\begin{ex}
    Khảo sát về cân nặng của các học sinh lớp 11D3 người ta được mẫu dữ liệu ghép nhóm như sau:
    \begin{center}
        \begin{tabular}{|c|c|c|c|c|c|c|}
            \hline
            Cân nặng & $[30;40)$ & $[40;50)$ & $[50;60)$ & $[60;70)$ & $[70;80)$ & $[80;90)$ \\
            \hline
            Số học sinh & $2$ & $10$ & $16$ & $8$ & $2$ & $2$ \\
            \hline
        \end{tabular}
    \end{center}
    Tứ phân vị thứ ba của mẫu số liệu trên là
    \choice
    {$60$}
    {\True $62,5$}
    {$65$}
    {$70$}
    \loigiai{
        Cỡ mẫu $N = 40$. Vị trí của $Q_3$ là $\frac{3N}{4} = 30$.
        Tần số tích lũy đến nhóm 3 là $2 + 10 + 16 = 28$, đến nhóm 4 là $28 + 8 = 36$.
        Nhóm chứa $Q_3$ là $[60;70)$.
        $Q_3 = 60 + \frac{30 - 28}{8} \cdot 10 = 62,5$.
    }
\end{ex}
\begin{ex}
    Khảo sát thời gian tập thể dục của một số học sinh khối 11 thu được mẫu số liệu ghép nhóm sau:
    \begin{center}
        \begin{tabular}{|c|c|c|c|c|c|}
            \hline
            Thời gian (phút) & $[0;20)$ & $[20;40)$ & $[40;60)$ & $[60;80)$ & $[80;100)$ \\
            \hline
            Số học sinh & $5$ & $9$ & $12$ & $10$ & $6$ \\
            \hline
        \end{tabular}
    \end{center}
    Nhóm chứa trung vị của mẫu số liệu trên là
    \choice
    {\True $[40;60)$}
    {$[20;40)$}
    {$[60;80)$}
    {$[80;100)$}
    \loigiai{
        Cỡ mẫu $N = 5 + 9 + 12 + 10 + 6 = 42$. Vị trí trung vị là $\frac{N}{2} = 21$.
        Tần số tích lũy đến nhóm 2 là $5 + 9 = 14$, đến nhóm 3 là $14 + 12 = 26$.
        Nhóm chứa trung vị là $[40;60)$.
    }
\end{ex}
\begin{ex}
    Khảo sát thời gian tập thể dục của một số học sinh khối 11 thu được mẫu số liệu ghép nhóm sau:
    \begin{center}
        \begin{tabular}{|c|c|c|c|c|c|}
            \hline
            Thời gian (phút) & $[0;20)$ & $[20;40)$ & $[40;60)$ & $[60;80)$ & $[80;100)$ \\
            \hline
            Số học sinh & $5$ & $9$ & $12$ & $10$ & $6$ \\
            \hline
        \end{tabular}
    \end{center}
    Nhóm chứa tứ phân vị thứ nhất của mẫu số liệu trên là
    \choice
    {$[40;60)$}
    {\True $[20;40)$}
    {$[60;80)$}
    {$[80;100)$}
    \loigiai{
        Cỡ mẫu $N = 42$. Vị trí tứ phân vị thứ nhất là $\frac{N}{4} = 10,5$.
        Tần số tích lũy đến nhóm 1 là $5$, đến nhóm 2 là $5 + 9 = 14$.
        Nhóm chứa tứ phân vị thứ nhất là $[20;40)$.
    }
\end{ex}

\begin{ex}
    Kết quả khảo sát năng suất (đơn vị: tấn/ha) của một số thửa ruộng được minh họa ở biểu đồ sau:\\
    \centerline{
        \begin{tikzpicture}[>=stealth, x=1.2cm, y=.6cm]
            % Vẽ các đường lưới ngang màu xanh nhạt
            \foreach \y in {1, 2, 3, 4, 5, 6} {
                \draw[help lines, blue!40] (0, \y) -- (7.5, \y);
                \node[left=0.1cm] at (0, \y) {\y}; % Nhãn trục tung
            }
            % Vẽ các cột histogram 
            % (Bắt đầu từ x=1 để cách trục tung một khoảng giống hình)
            \begin{scope}[fill=cyan!50, draw=black]
                \foreach \x/\y in {1/3, 2/4, 3/6, 4/5, 5/5, 6/2} {
                    \filldraw (\x, 0) rectangle (\x+1, \y);
                }
            \end{scope}
            % Vẽ trục tọa độ
            \draw[->, thick] (0, 0) -- (8, 0) node[below, align=center] {Năng suất\\(tấn/ha)};
            \draw[->, thick] (0, 0) -- (0, 6.8) node[left] {Số thửa ruộng};
            \node[below left] at (0, 0) {$0$};
            % Các vạch chia và nhãn trên trục hoành
            \foreach \x/\val in {1/5{,}5, 2/5{,}7, 3/5{,}9, 4/6{,}1, 5/6{,}3, 6/6{,}5, 7/6{,}7} {
                \draw (\x, 0.1) -- (\x, -0.1) node[below=0.15cm] {$\val$};
            }
        \end{tikzpicture}
    }
    Khoảng biến thiên $R$ của mẫu số liệu trên là
    \choice
    {$R = 4$}
    {\True $R = 1,2$}
    {$R = 6$}
    {$R = 6,7$}
%    \loigiai{
%        Dựa vào biểu đồ, các nhóm số liệu trải dài từ giá trị nhỏ nhất là $5,5$ đến giá trị lớn nhất là $6,7$.
%        Khoảng biến thiên của mẫu số liệu là: $R = 6,7 - 5,5 = 1,2$.
%    }
\end{ex}
\begin{ex}
    Khảo sát thời gian xem ti vi trong một ngày của một số học sinh khối 11 thu được mẫu số liệu ghép nhóm sau:\\
    \centerline{
        \begin{tabular}{|l|c|c|c|c|c|}
            \hline
            Thời gian (phút) & $[0; 20)$ & $[20; 40)$ & $[40; 60)$ & $[60; 80)$ & $[80; 100)$ \\
            \hline
            Số học sinh & $5$ & $9$ & $12$ & $10$ & $6$ \\
            \hline
        \end{tabular}
    }
    Khoảng biến thiên $R$ của mẫu số liệu ghép nhóm trên là
    \choice
    {$R = 20$}
    {\True $R = 100$}
    {$R = 50$}
    {$R = 60$}
%    \loigiai{
%        Đầu mút phải của nhóm cuối cùng là $100$, đầu mút trái của nhóm đầu tiên là $0$.
%        Khoảng biến thiên: $R = 100 - 0 = 100$.
%    }
\end{ex}
\begin{ex}
    Thống kê thời gian sử dụng mạng xã hội trong ngày của các bạn học sinh tổ 1 và tổ 2 lớp 12A thu được bảng sau:\\
    \centerline{
        \begin{tabular}{|l|c|c|c|c|}
            \hline
            Thời gian sử dụng (phút) & $[0; 10)$ & $[10; 30)$ & $[30; 60)$ & $[60; 90)$ \\
            \hline
            Số học sinh Tổ 1 & $2$ & $4$ & $3$ & $1$ \\
            \hline
            Số học sinh Tổ 2 & $5$ & $1$ & $3$ & $0$ \\
            \hline
        \end{tabular}
    }
    Tìm khoảng biến thiên $R_1, R_2$ cho thời gian sử dụng mạng xã hội của tổ 1 và tổ 2.
    \choice
    {$R_1 = 60; R_2 = 90$}
    {\True $R_1 = 90; R_2 = 60$}
    {$R_1 = 30; R_2 = 30$}
    {$R_1 = 10; R_2 = 30$}
%    \loigiai{
%        - Đối với tổ 1: Nhóm chứa dữ liệu lớn nhất là $[60; 90)$ và nhỏ nhất là $[0; 10)$. Vậy $R_1 = 90 - 0 = 90$.
%        - Đối với tổ 2: Nhóm chứa dữ liệu lớn nhất là $[30; 60)$ (nhóm $[60; 90)$ có tần số bằng 0) và nhỏ nhất là $[0; 10)$. Vậy $R_2 = 60 - 0 = 60$.
%    }
\end{ex}
\begin{ex}
    Tìm hiểu thời gian hoàn thành một bài tập (đơn vị: phút) của một số học sinh thu được kết quả sau:
    \begin{center}
        \begin{tabular}{|l|c|c|c|c|c|}
            \hline
            Thời gian (phút) & $[0; 4]$ & $[4; 8)$ & $[8; 12)$ & $[12; 16)$ & $[16; 20)$ \\
            \hline
            Số học sinh & $2$ & $4$ & $7$ & $4$ & $3$ \\
            \hline
        \end{tabular}
    \end{center}
    Tìm khoảng biến thiên cho mẫu số liệu ghép nhóm trên
    \choice
    {\True $20$}
    {$15$}
    {$16$}
    {$4$}
%    \loigiai{
%        Đầu mút lớn nhất của nhóm cuối cùng là $20$, đầu mút nhỏ nhất của nhóm đầu tiên là $0$.
%        Khoảng biến thiên: $R = 20 - 0 = 20$.
%    }
\end{ex}
\begin{ex}
    Kết quả khảo sát năng suất (đơn vị: tấn/ha) của một số thửa ruộng được minh họa ở biểu đồ sau\\
    \centerline{
        \begin{tikzpicture}[>=stealth, x=1.2cm, y=.6cm]
            % Vẽ các đường lưới ngang màu xanh nhạt
            \foreach \y in {1, 2, 3, 4, 5, 6} {
                \draw[help lines, blue!40] (0, \y) -- (7.5, \y);
                \node[left=0.1cm] at (0, \y) {\y}; % Nhãn trục tung
            }
            % Vẽ các cột histogram 
            % (Bắt đầu từ x=1 để cách trục tung một khoảng giống hình)
            \begin{scope}[fill=cyan!50, draw=black]
                \foreach \x/\y in {1/3, 2/4, 3/6, 4/5, 5/5, 6/2} {
                    \filldraw (\x, 0) rectangle (\x+1, \y);
                }
            \end{scope}
            % Vẽ trục tọa độ
            \draw[->, thick] (0, 0) -- (8, 0) node[below, align=center] {Năng suất\\(tấn/ha)};
            \draw[->, thick] (0, 0) -- (0, 6.8) node[left] {Số thửa ruộng};
            \node[below left] at (0, 0) {$0$};
            % Các vạch chia và nhãn trên trục hoành
            \foreach \x/\val in {1/5{,}5, 2/5{,}7, 3/5{,}9, 4/6{,}1, 5/6{,}3, 6/6{,}5, 7/6{,}7} {
                \draw (\x, 0.1) -- (\x, -0.1) node[below=0.15cm] {$\val$};
            }
        \end{tikzpicture}
    }
    Khoảng tứ phân vị của mẫu số liệu ghép nhóm đã cho là (làm tròn đến hàng phần trăm).
    \choice
    {$0,4576$}
    {$0,4765$}
    {\True $0,4675$}
    {$0,4657$}
    \loigiai{
        Cỡ mẫu $n = 3 + 4 + 6 + 5 + 5 + 2 = 25$. Độ dài nhóm ghép $h=0,2$.
        - $Q_1$ nằm ở vị trí thứ $\frac{25}{4} = 6,25$, thuộc nhóm $[5,7; 5,9)$: $Q_1 = 5,7 + \frac{6,25 - 3}{4} \cdot 0,2 = 5,8625$.
        - $Q_3$ nằm ở vị trí thứ $\frac{3 \cdot 25}{4} = 18,75$, thuộc nhóm $[6,3; 6,5)$: $Q_3 = 6,3 + \frac{18,75 - 18}{5} \cdot 0,2 = 6,33$.
        Khoảng tứ phân vị: $\Delta_Q = Q_3 - Q_1 = 6,33 - 5,8625 = 0,4675$.
    }
\end{ex}
\begin{ex}
    Cho mẫu số liệu ghép nhóm về độ tuổi của dân cư của khu phố $A$ như sau\\
    \centerline{
        \begin{tabular}{|c|c|c|c|c|c|c|}
            \hline
            Nhóm & $[20; 30)$ & $[30; 40)$ & $[40; 50)$ & $[50; 60)$ & $[60; 70)$ & $[70; 80)$ \\
            \hline
            Số người & $24$ & $26$ & $20$ & $15$ & $11$ & $4$ \\
            \hline
        \end{tabular}
    }
    Khoảng tứ phân vị của mẫu số liệu ghép nhóm đã cho là (làm tròn đến hàng phần trăm).
    \choice
    {$30,38$}
    {$63,33$}
    {\True $22,95$}
    {$22,94$}
%    \loigiai{
%        Cỡ mẫu $n = 100$.
%        - $Q_1$ ở vị trí thứ $\frac{100}{4}=25$, thuộc nhóm $[30; 40)$: $Q_1 = 30 + \frac{25 - 24}{26} \cdot 10 \approx 30,38$.
%        - $Q_3$ ở vị trí thứ $\frac{3 \cdot 100}{4}=75$, thuộc nhóm $[50; 60)$: $Q_3 = 50 + \frac{75 - 70}{15} \cdot 10 \approx 53,33$.
%        Khoảng tứ phân vị: $\Delta_Q = 53,33 - 30,38 = 22,95$.
%    }
\dotlineEX{1}
\end{ex}
\begin{ex}
    Lương tháng của một số nhân viên một văn phòng được ghi lại như sau (đơn vị: triệu đồng):\\
    \centerline{
        \begin{tabular}{|l|c|c|c|c|}
            \hline
            Lương tháng (triệu đồng) & $[6; 8)$ & $[8; 10)$ & $[10; 12)$ & $[12; 14)$ \\
            \hline
            Số nhân viên & $3$ & $6$ & $8$ & $7$ \\
            \hline
        \end{tabular}
    }
    Tìm khoảng tứ phân vị của dãy số liệu trên.
    \choice
    {$1,75$}
    {\True $3,3$}
    {$3,25$}
    {$2,3$}
%    \loigiai{
%        Cỡ mẫu $n = 24$.
%        - $Q_1$ nằm ở vị trí $\frac{24}{4} = 6$, thuộc nhóm $[8; 10)$: $Q_1 = 8 + \frac{6-3}{6} \cdot 2 = 9$.
%        - $Q_3$ nằm ở vị trí $\frac{3 \cdot 24}{4} = 18$, thuộc nhóm $[12; 14)$: $Q_3 = 12 + \frac{18-17}{7} \cdot 2 \approx 12,285$.
%        Khoảng tứ phân vị $\Delta_Q \approx 12,285 - 9 = 3,285 \approx 3,3$.
%    }
\dotlineEX{1}
\end{ex}
\begin{ex}
    Số lượng học sinh trên lớp đăng ký tham gia hoạt động Hoa phượng đỏ ở một trường THPT trên địa bàn TP.HCM được cho ở bảng sau:\\
    \centerline{
        \begin{tabular}{|l|c|c|c|c|}
            \hline
            Điểm số & $[6; 10]$ & $[11; 15]$ & $[16; 20]$ & $[21; 25]$ \\
            \hline
            Số học sinh & $4$ & $8$ & $2$ & $6$ \\
            \hline
        \end{tabular}
    }
    Giá trị nào sau đây là giá trị ngoại lệ của mẫu số liệu trên
    \choice
    {\True $38$}
    {$9$}
    {$15$}
    {$10$}
    \loigiai{
        Chuyển về nhóm nửa khoảng: $[5,5; 10,5), [10,5; 15,5), [15,5; 20,5), [20,5; 25,5)$. $n = 20$.
        - $Q_1$ vị trí $5$ thuộc $[10,5; 15,5)$: $Q_1 = 10,5 + \frac{5-4}{8} \cdot 5 = 11,125$.
        - $Q_3$ vị trí $15$ thuộc $[20,5; 25,5)$: $Q_3 = 20,5 + \frac{15-14}{6} \cdot 5 = 21,333$.
        $\Delta_Q = 10,208$. Giới hạn trên để không là ngoại lệ: $Q_3 + 1,5\Delta_Q \approx 36,6$. Vì $38 > 36,6$ nên $38$ là ngoại lệ.
    }
\end{ex}
\begin{ex}
    Các bạn học sinh lớp 12A5 trả lời 40 câu hỏi trong một bài kiểm tra. Kết quả số câu trả lời đúng được thống kê ở bảng sau.
    \begin{center}
        \begin{tabular}{|l|c|c|c|c|c|}
            \hline
            Số câu trả lời đúng & $[16; 21)$ & $[21; 26)$ & $[26; 31)$ & $[31; 36)$ & $[36; 41)$ \\
            \hline
            Số học sinh & $4$ & $8$ & $8$ & $16$ & $4$ \\
            \hline
        \end{tabular}
    \end{center}
    Khoảng tứ phân vị của mẫu số liệu là
    \choice
    {\True $9,375$}
    {$8,625$}
    {$10,15$}
    {$7,5$}
    \loigiai{
        Cỡ mẫu $n = 40$.
        - $Q_1$ vị trí $10$ thuộc $[21; 26)$: $Q_1 = 21 + \frac{10-4}{8} \cdot 5 = 24,75$.
        - $Q_3$ vị trí $30$ thuộc $[31; 36)$: $Q_3 = 31 + \frac{30-20}{16} \cdot 5 = 34,125$.
        Khoảng tứ phân vị $\Delta_Q = 34,125 - 24,75 = 9,375$.
    }
\end{ex}
\begin{ex}
    Cho mẫu số liệu ghép nhóm về chiều cao của 25 cây cam giống như sau:
    \begin{center}
        \begin{tabular}{|l|c|c|c|c|c|}
            \hline
            Chiều cao (cm) & $[0; 10)$ & $[10; 20)$ & $[20; 30)$ & $[30; 40)$ & $[40; 50)$ \\
            \hline
            Số cây & $4$ & $6$ & $7$ & $5$ & $3$ \\
            \hline
        \end{tabular}
    \end{center}
    Từ một mẫu số liệu về chiều cao của cây xoài giống người ta tính được khoảng tứ phân vị bằng $13,94$. Đối với các cây cam giống và xoài giống được khảo sát ở trên, khẳng định nào sau đây đúng
    \choice
    {Chiều cao của các cây xoài giống phân tán hơn}
    {\True Chiều cao của các cây cam giống phân tán hơn}
    {Các cây cam và xoài giống có chiều cao phân tán như nhau}
    {Không so sánh được độ phân tán của các cây cam giống và xoài giống được khảo sát}
    \loigiai{
        Đối với cây cam: $n=25$.
        - $Q_1$ vị trí $6,25$ thuộc $[10; 20)$: $Q_1 = 10 + \frac{6,25-4}{6} \cdot 10 = 13,75$.
        - $Q_3$ vị trí $18,75$ thuộc $[30; 40)$: $Q_3 = 30 + \frac{18,75-17}{5} \cdot 10 = 33,5$.
        $\Delta_Q \text{ (cam)} = 33,5 - 13,75 = 19,75$.
        Vì $19,75 > 13,94$ nên chiều cao cây cam giống phân tán hơn.
    }
    \dotlineEX{1}
\end{ex}
\begin{ex}
    Một công ty xây dựng khảo sát khách hàng xem họ có nhu cầu mua nhà ở mức giá nào. Kết quả khảo sát được lưu lại ở bảng sau:
    \begin{center}
        \begin{tabular}{|c|c|c|c|c|c|}
            \hline
            Mức giá (triệu đồng/m$^2$) & $[10;14)$ & $[14;18)$ & $[18;22)$ & $[22;26)$ & $[26;30)$ \\
            \hline
            Số khách hàng & $54$ & $78$ & $120$ & $45$ & $12$ \\
            \hline
        \end{tabular}
    \end{center}
    \choiceTF
    {\True Độ dài của nhóm $[10;14)$ bằng $4$}
    {\True Giá trị đại diện của nhóm $[10;14)$ là $12$}
    {Số trung bình của mẫu số liệu ghép nhóm bằng $19,52$}
    {\True Mốt của mẫu số liệu ghép nhóm là $19,44$}
    \loigiai{
        \begin{itemize}
            \item Ý a đúng: Độ dài nhóm $= 14 - 10 = 4$.
            \item Ý b đúng: Giá trị đại diện $= \frac{10 + 14}{2} = 12$.
            \item Ý c sai: Cỡ mẫu $N = 309$. $\bar{x} = \frac{12 \cdot 54 + 16 \cdot 78 + 20 \cdot 120 + 24 \cdot 45 + 28 \cdot 12}{309} \approx 18,49 \neq 19,52$.
            \item Ý d đúng: Nhóm chứa mốt là $[18;22)$. $M_o = 18 + \frac{120 - 78}{(120 - 78) + (120 - 45)} \cdot 4 \approx 19,44$.
        \end{itemize}
    }
    \dotlineEX{1}
\end{ex}
\begin{ex}
    Phỏng vấn một số học sinh khối 11 về thời gian (giờ) ngủ của một buổi tối, thu được bảng số liệu sau:
    \begin{center}
        \begin{tabular}{|c|c|c|}
            \hline
            Thời gian & Số học sinh nam & Số học sinh nữ \\
            \hline
            $[4;5)$ & $8$ & $4$ \\
            \hline
            $[5;6)$ & $10$ & $6$ \\
            \hline
            $[6;7)$ & $13$ & $12$ \\
            \hline
            $[7;8)$ & $15$ & $18$ \\
            \hline
            $[8;9)$ & $7$ & $8$ \\
            \hline
        \end{tabular}
    \end{center}
    \choiceTF
    {\True Độ dài của nhóm bằng $1$}
    {Thời gian ngủ trung bình của các bạn học sinh nam nhiều hơn các bạn học sinh nữ}
    {\True Phần lớn học sinh được khảo sát trong khối 11 ngủ nhiều hơn $6,5$ giờ}
    {$75\%$ học sinh được khảo sát trong khối 11 ngủ ít nhất $5,5$ giờ}
    \loigiai{
        Tổng số học sinh Nam $= 53$, Nữ $= 48$, cả khối $= 101$.
        \begin{itemize}
            \item Ý a đúng: Các nhóm đều có độ dài $5-4=1$.
            \item Ý b sai: $\bar{x}_{\text{Nam}} \approx 6,56$ giờ, $\bar{x}_{\text{Nữ}} \approx 6,92$ giờ. Nam ít hơn Nữ.
            \item Ý c đúng: Học sinh ngủ $> 6,5$ giờ tương đương từ nửa sau của nhóm $[6;7)$ trở đi. Số lượng ước tính là $\frac{25}{2} + 33 + 15 = 60,5$ học sinh, chiếm khoảng $60\%$ (lớn hơn $50\%$).
            \item Ý d sai: Theo tính toán tứ phân vị thứ nhất $Q_1 \approx 5,82$. Điều này có nghĩa $75\%$ học sinh ngủ ít nhất $5,82$ giờ, giá trị này khác với mốc $5,5$ giờ được đưa ra để đánh lừa. (Thực tế có khoảng $80,2\%$ học sinh ngủ $\ge 5,5$ giờ).
        \end{itemize}
    }
\end{ex}
\begin{ex}
    Số lượng người đi xem bộ phim Lật Mặt 7 của Lý Hải theo độ tuổi trong một rạp chiếu phim (sau 1h đầu công chiếu) được ghi lại theo bảng số liệu sau:
    \begin{center}
        \begin{tabular}{|c|c|c|c|c|c|}
            \hline
            Độ tuổi & $[10;20)$ & $[20;30)$ & $[30;40)$ & $[40;50)$ & $[50;60)$ \\
            \hline
            Số người & $19$ & $28$ & $16$ & $7$ & $4$ \\
            \hline
        \end{tabular}
    \end{center}
    \choiceTF
    {\True Cỡ mẫu của mẫu số liệu là $74$}
    {\True Độ tuổi được dự báo là ít xem phim Lật Mặt 7 nhất là thuộc nhóm $[50;60)$}
    {Trung vị của mẫu số liệu là $26,42$ (kết quả làm tròn đến hàng phần trăm)}
    {Độ tuổi được dự báo là thích xem phim Lật Mặt 7 nhiều nhất là $29$ tuổi}
    \loigiai{
        \begin{itemize}
            \item Ý a đúng: $N = 19 + 28 + 16 + 7 + 4 = 74$.
            \item Ý b đúng: Nhóm $[50;60)$ có tần số thấp nhất là $4$.
            \item Ý c sai: Nhóm chứa trung vị là $[20;30)$. $M_e = 20 + \frac{37 - 19}{28} \cdot 10 \approx 26,43$ (làm tròn phần trăm).
            \item Ý d sai: Độ tuổi dự báo xem nhiều nhất là Mốt. $M_o = 20 + \frac{28 - 19}{(28 - 19) + (28 - 16)} \cdot 10 \approx 24,28$ tuổi, không phải $29$.
        \end{itemize}
    }
\end{ex}

\begin{ex}
    Cho bảng mẫu số liệu ghép nhóm về mức lương của hai công ty $A$ và $B$ (đơn vị: triệu đồng) như sau:
    \begin{center}
        \begin{tabular}{|c|c|c|c|c|c|c|c|}
            \hline
            Mức lương & $[5; 9)$ & $[9; 13)$ & $[13; 17)$ & $[17; 21)$ & $[21; 25)$ & $[25; 29)$ & $[29; 34)$ \\
            \hline
            Số người công ty A & $10$ & $37$ & $20$ & $13$ & $13$ & $5$ & $2$ \\
            \hline
            Số người công ty B & $16$ & $38$ & $27$ & $11$ & $10$ & $3$ & $0$ \\
            \hline
        \end{tabular}
    \end{center}
    \choiceTF
    {Khoảng biến thiên cho mức lương của hai công ty là $29$}
    {\True Khoảng biến thiên cho mức lương của công ty $B$ là $24$}
    {\True Mức lương của công ty $A$ biến động nhiều hơn mức lương của công ty $B$}
    {\True Nhóm chứa tứ phân vị thứ ba của công ty $A$ là nhóm $[17; 21)$}
    \loigiai{
        - Khoảng biến thiên công ty A là $34-5 = 29$, công ty B là $29-5 = 24$. Ý a) sai, ý b) đúng.
        - Vì khoảng biến thiên của A lớn hơn B nên mức lương công ty A biến động nhiều hơn. Ý c) đúng.
        - Tổng số người công ty A là $100$. $Q_3$ nằm ở vị trí $75$. Tần số tích lũy: $10, 47, 67, 80$. Do đó $Q_3$ nằm ở nhóm $[17; 21)$. Ý d) đúng.
    }
\end{ex}
\begin{ex}
    Thời gian (phút) truy cập Internet mỗi buổi tối của một số học sinh được cho ở bảng sau
    \begin{center}
        \begin{tabular}{|c|c|c|c|c|c|}
            \hline
            Thời gian (phút) & $[9,5; 12,5)$ & $[12,5; 15,5)$ & $[15,5; 18,5)$ & $[18,5; 21,5)$ & $[21,5; 24,5)$ \\
            \hline
            Số học sinh & $4$ & $12$ & $14$ & $23$ & $3$ \\
            \hline
        \end{tabular}
    \end{center}
    \choiceTF
    {\True Khoảng biến thiên của mẫu số liệu là $15$}
    {Nhóm chứa tứ phân vị thứ ba là $[15,5; 18,5)$}
    {\True Tứ phân vị thứ nhất là $Q_1 = 15$}
    {\True Khoảng tứ phân vị của mẫu số liệu ghép nhóm bé hơn $6$}
%    \loigiai{
%        - Khoảng biến thiên $R = 24,5 - 9,5 = 15$. Ý a) đúng.
%        - Cỡ mẫu $n = 56$. $Q_3$ nằm ở vị trí $42$, thuộc nhóm $[18,5; 21,5)$. Ý b) sai.
%        - $Q_1$ nằm ở vị trí $14$, thuộc nhóm $[12,5; 15,5)$. $Q_1 = 12,5 + \frac{14-4}{12} \cdot 3 = 15$. Ý c) đúng.
%        - $Q_3 = 18,5 + \frac{42-30}{23} \cdot 3 \approx 20,06$. $\Delta_Q = 20,06 - 15 = 5,06 < 6$. Ý d) đúng.
%    }
\end{ex}

\begin{ex}
    Bảng tần số ghép nhóm dưới đây thống kê số giờ ngủ buổi tối của các học sinh lớp 12A1 và 12A2:
    \begin{center}
        \begin{tabular}{|l|c|c|c|c|c|}
            \hline
            Thời gian & $[4; 5)$ & $[5; 6)$ & $[6; 7)$ & $[7; 8)$ & $[8; 9)$ \\
            \hline
            Số học sinh nam & $6$ & $10$ & $13$ & $9$ & $7$ \\
            \hline
            Số học sinh nữ & $4$ & $8$ & $10$ & $11$ & $8$ \\
            \hline
        \end{tabular}
    \end{center}
    \choiceTF
    {\True Khoảng biến thiên của mẫu số liệu ghép nhóm về thời gian ngủ của các bạn nam là $5$}
    {Khoảng tứ phân vị của mẫu số liệu ghép nhóm về thời gian ngủ của các bạn nam là $2,09$}
    {\True Khoảng tứ phân vị của mẫu số liệu ghép nhóm về thời gian ngủ của các bạn nữ trong khoảng $(2; 3)$}
    {\True Học sinh nam có thời gian ngủ đồng đều hơn.}
%    \loigiai{
%        - Khoảng biến thiên của nam $R = 9-4 = 5$. Ý a) đúng.
%        - Với nam ($n=45$): $\Delta_Q \approx 7,53 - 5,52 = 2,01 \neq 2,09$. Ý b) sai.
%        - Với nữ ($n=41$): $\Delta_Q \approx 7,80 - 5,78 = 2,02$. Giá trị này thuộc khoảng $(2;3)$. Ý c) đúng.
%        - Vì $\Delta_Q$ nam (khoảng $2,00$) nhỏ hơn $\Delta_Q$ nữ (khoảng $2,01$) nên dữ liệu nam đồng đều hơn. Ý d) đúng.
%    }
\end{ex}
\begin{ex}
    Biểu đồ dưới đây thống kê thời gian tập thể dục buổi sáng mỗi ngày trong tháng 9/2022 của bác Bình và bác An.
    \begin{center}
        \begin{tikzpicture}[>=stealth, scale=.7]
            % Định nghĩa màu sắc giống trong hình
            \definecolor{colorBinh}{RGB}{68, 84, 143} % Xanh đậm (Bác Bình)
            \definecolor{colorAn}{RGB}{255, 153, 35} % Cam (Bác An)
            % Tiêu đề biểu đồ
            \node[font=\large\bfseries] at (5.5, 7) {THỜI GIAN TẬP THỂ DỤC};
            % Vẽ các đường lưới ngang và nhãn trục tung
            \foreach \y/\val in {1/5, 2/10, 3/15, 4/20, 5/25, 6/30} {
                \draw[help lines, gray!50] (0, \y) -- (11.5, \y);
                \node[left=0.1cm] at (0, \y) {\val};
            }
            % Gốc tọa độ
            \node[left=0.1cm] at (0, 0) {$0$};
            % Vẽ các cột bằng vòng lặp
            % Cú pháp: \x (vị trí tâm nhóm), \label (nhãn trục hoành), \valB (Bình), \valA (An)
            \foreach \x/\label/\valB/\valA in {
                1.5/{[15; 20)}/5/0,
                3.5/{[20; 25)}/12/25,
                5.5/{[25; 30)}/8/5,
                7.5/{[30; 35)}/3/0,
                9.5/{[35; 40)}/2/0
            } {
                % In nhãn dưới trục hoành
                \node[below=0.2cm] at (\x, 0) {$\label$};
                % Cột của Bác Bình (vẽ lệch sang trái)
                \fill[colorBinh] (\x-0.8, 0) rectangle (\x, \valB/5);
                \node[above] at (\x-0.4, \valB/5) {\valB};
                % Cột của Bác An (vẽ lệch sang phải)
                \fill[colorAn] (\x, 0) rectangle (\x+0.8, \valA/5);
                \node[above] at (\x+0.4, \valA/5) {\valA};
            }
            % Vẽ trục tung và trục hoành (vẽ sau cùng để đè lên các chi tiết khác)
            \draw[->, thick] (0, 0) -- (0, 7) node[left, align=center] {Số\\ngày};
            \draw[->, thick] (0, 0) -- (12, 0) node[below=0.1cm, align=center] {Thời gian\\(phút)};
            % Chú giải (Legend) ở bên phải biểu đồ
            \begin{scope}[shift={(11.5, 3.5)}]
                \fill[colorBinh] (0, 1) rectangle (0.3, 1.3);
                \node[right] at (0.4, 1.15) {Bác Bình};
                \fill[colorAn] (0, 0) rectangle (0.3, 0.3);
                \node[right] at (0.4, 0.15) {Bác An};
            \end{scope}
        \end{tikzpicture}
    \end{center}
    \choiceTF
    {\True Khoảng biến thiên của mẫu số liệu ghép nhóm về thời gian tập thể dục buổi sáng của bác Bình là $25$ (phút).}
    {Khoảng tứ phân vị của mẫu số liệu ghép nhóm về thời gian tập thể dục buổi sáng của bác An là: $\Delta_Q = 2$}
    {\True Tứ phân vị thứ ba của mẫu số liệu ghép nhóm về thời gian tập thể dục buổi sáng của bác Bình là: $Q_3' = \frac{455}{16}$}
    {Khoảng tứ phân vị của mẫu số liệu ghép nhóm về thời gian tập thể dục buổi sáng mỗi ngày của bác An lớn hơn bác Bình}
    \loigiai{
        - Bác Bình ($n=30$): Khoảng biến thiên là $40 - 15 = 25$. Ý a) đúng.
        - Bác An ($n=30$): $Q_1 = 20 + \frac{7,5}{25} \cdot 5 = 21,5$; $Q_3 = 20 + \frac{22,5}{25} \cdot 5 = 24,5$. $\Delta_Q = 3 \neq 2$. Ý b) sai.
        - Bác Bình ($n=30$): $Q_3$ nằm ở nhóm $[25; 30)$: $Q_3' = 25 + \frac{22,5 - 17}{8} \cdot 5 = \frac{455}{16}$. Ý c) đúng.
        - $\Delta_Q$ bác Bình bằng $\frac{455}{16} - \left(20 + \frac{7,5-5}{12} \cdot 5\right) \approx 7,4$, lớn hơn của bác An ($3$). Ý d) sai.
    }
\end{ex}


\begin{ex}
    Bảng thống kê nhiệt độ (C) cao nhất trong các ngày của tháng 7 tại tỉnh Hà Tĩnh trong năm 2024 cho ở bảng sau:
    \begin{center}
        \begin{tabular}{|c|c|c|c|c|c|c|}
            \hline
            Nhiệt độ (C) & $[28; 30)$ & $[30; 32)$ & $[32; 34)$ & $[34; 36)$ & $[36; 38)$ & $[38; 40)$ \\
            \hline
            Số ngày & $1$ & $1$ & $6$ & $12$ & $9$ & $2$ \\
            \hline
        \end{tabular}
    \end{center}
    Ngày nóng nhất và ngày mát nhất trong tháng 7 chênh nhau bao nhiêu độ (C)?
    \shortans{$12$}
    \loigiai{
        Khoảng biến thiên của mẫu số liệu là $R = 40 - 28 = 12$.
        Vậy ngày nóng nhất và ngày mát nhất trong tháng 7 chênh nhau tối đa $12^\circ\text{C}$.
    }
\end{ex}
\begin{ex}
    Khoảng tứ phân vị của mẫu số liệu ghép nhóm được cho ở bảng sau là bao nhiêu? (làm tròn kết quả đến hai chữ số thập phân).
    \begin{center}
        \begin{tabular}{|c|c|c|c|c|}
            \hline
            Nhóm & $[25; 30)$ & $[30; 35)$ & $[35; 40)$ & $[40; 45)$ \\
            \hline
            Tần số & $2$ & $17$ & $10$ & $25$ \\
            \hline
        \end{tabular}
    \end{center}
    \shortans{$8,92$}
    \loigiai{
        Cỡ mẫu $n = 54$.
        - Vị trí của $Q_1$ là $\frac{54}{4} = 13,5$. Nhóm chứa $Q_1$ là $[30; 35)$.
        $Q_1 = 30 + \frac{13,5 - 2}{17} \cdot 5 \approx 33,38$.
        - Vị trí của $Q_3$ là $\frac{3 \cdot 54}{4} = 40,5$. Nhóm chứa $Q_3$ là $[40; 45)$.
        $Q_3 = 40 + \frac{40,5 - 29}{25} \cdot 5 = 42,3$.
        - Khoảng tứ phân vị $\Delta_Q = Q_3 - Q_1 = 42,3 - 33,38 = 8,92$.
    }
\end{ex}
\begin{ex}
    Cho mẫu số liệu về thời gian (phút) đi từ nhà đến trường của các học sinh trong một lớp 12 của một trường như sau:
    \begin{center}
        \begin{tabular}{|c|c|c|c|c|c|c|}
            \hline
            Thời gian & $[0; 5)$ & $[5; 10)$ & $[10; 15)$ & $[15; 20)$ & $[20; 25)$ & $[25; 30)$ \\
            \hline
            Số học sinh & $7$ & $12$ & $7$ & $5$ & $3$ & $2$ \\
            \hline
        \end{tabular}
    \end{center}
    Hãy tìm khoảng tứ phân vị của mẫu số liệu ghép nhóm trên. (kết quả làm tròn đến hàng phần mười)
    \shortans{$8,9$}
    \loigiai{
        Cỡ mẫu $n = 36$.
        - Vị trí của $Q_1$ là $\frac{36}{4} = 9$. Nhóm chứa $Q_1$ là $[5; 10)$.
        $Q_1 = 5 + \frac{9 - 7}{12} \cdot 5 = 5 + \frac{5}{6} \approx 5,83$.
        - Vị trí của $Q_3$ là $\frac{3 \cdot 36}{4} = 27$. Nhóm chứa $Q_3$ là $[15; 20)$.
        $Q_3 = 15 + \frac{27 - 26}{5} \cdot 5 = 16$.
        - Khoảng tứ phân vị $\Delta_Q = Q_3 - Q_1 = 16 - \frac{35}{6} \approx 10,17 \approx 10,2$.
    }
\end{ex}
\begin{ex}
    Bảng sau thống kê cân nặng của 50 quả xoài Thanh Ca được lựa chọn ngẫu nhiên sau khi thu hoạch ở một nông trường.
    \begin{center}
        \begin{tabular}{|c|c|c|c|c|c|}
            \hline
            Cân nặng (g) & $[250; 290)$ & $[290; 330)$ & $[330; 370)$ & $[370; 410)$ & $[410; 450)$ \\
            \hline
            Số quả xoài & $3$ & $13$ & $18$ & $11$ & $5$ \\
            \hline
        \end{tabular}
    \end{center}
    Khoảng tứ phân vị của mẫu số liệu ghép nhóm trên bằng. (kết quả làm tròn đến hàng phần mười)
    \shortans{$64,6$}
    \loigiai{
        Cỡ mẫu $n = 50$.
        - Vị trí của $Q_1$ là $\frac{50}{4} = 12,5$. Nhóm chứa $Q_1$ là $[290; 330)$.
        $Q_1 = 290 + \frac{12,5 - 3}{13} \cdot 40 \approx 319,23$.
        - Vị trí của $Q_3$ là $\frac{3 \cdot 50}{4} = 37,5$. Nhóm chứa $Q_3$ là $[370; 410)$.
        $Q_3 = 370 + \frac{37,5 - 34}{11} \cdot 40 \approx 382,73$.
        - Khoảng tứ phân vị $\Delta_Q = Q_3 - Q_1 \approx 382,73 - 319,23 = 63,5$.
    }
\end{ex}

\begin{ex}
    Người ta ghi lại tuổi thọ của một số con ong cho kết quả như sau:
    \begin{center}
        \begin{tabular}{|c|c|c|c|c|c|}
            \hline
            Tuổi thọ (ngày) & $[0; 20)$ & $[20; 40)$ & $[40; 60)$ & $[60; 80)$ & $[80; 100)$ \\
            \hline
            Số lượng & $5$ & $12$ & $23$ & $31$ & $29$ \\
            \hline
        \end{tabular}
    \end{center}
    Tính khoảng tứ phân vị của mẫu số liệu. (kết quả làm tròn đến hàng phần mười)
    \shortans{$38,5$}
    \loigiai{
        Cỡ mẫu $n = 100$.
        - Vị trí của $Q_1$ là $\frac{100}{4} = 25$. Nhóm chứa $Q_1$ là $[40; 60)$.
        $Q_1 = 40 + \frac{25 - 17}{23} \cdot 20 \approx 46,96$.
        - Vị trí của $Q_3$ là $\frac{3 \cdot 100}{4} = 75$. Nhóm chứa $Q_3$ là $[80; 100)$.
        $Q_3 = 80 + \frac{75 - 71}{29} \cdot 20 \approx 82,76$.
        - Khoảng tứ phân vị $\Delta_Q = Q_3 - Q_1 \approx 82,76 - 46,96 = 35,8$.
    }
\end{ex}


\begin{ex}
    Thống kê về thời lượng mỗi trận đấu bi-a trong vòng tứ kết giải đấu European Open người ta được mẫu số liệu ghép nhóm như sau:
    \begin{center}
        \begin{tabular}{|c|c|c|c|c|c|}
            \hline
            Thời gian & $[9,5;12,5)$ & $[12,5;15,5)$ & $[15,5;18,5)$ & $[18,5;21,5)$ & $[21,5;24,5)$ \\
            \hline
            Số trận & $3$ & $12$ & $15$ & $24$ & $2$ \\
            \hline
        \end{tabular}
    \end{center}
    Số trung vị của mẫu số liệu trên là \textit{(kết quả làm tròn đến hàng phần mười)}.
    \shortans{18,1}
    \loigiai{
        Cỡ mẫu $N = 56$. Vị trí trung vị là $\frac{N}{2} = 28$.
        Tần số tích lũy đến nhóm 2 là $15$, đến nhóm 3 là $30$. Nhóm chứa trung vị là $[15,5;18,5)$.
        $M_e = 15,5 + \frac{28 - 15}{15} \cdot 3 = 18,1$.
    }
\end{ex}
\begin{ex}
    Kết quả điều tra chiều cao của học sinh (đơn vị: cm) trong trường THPT được biểu diễn ở biểu đồ Histogram, tương đương với bảng tần số sau:
    \begin{center}
        \begin{tabular}{|c|c|c|c|c|c|}
            \hline
            Chiều cao (cm) & $[160;163)$ & $[163;166)$ & $[166;169)$ & $[169;172)$ & $[172;175)$ \\
            \hline
            Số học sinh & $6$ & $12$ & $9$ & $5$ & $3$ \\
            \hline
        \end{tabular}
    \end{center}
    Tứ phân vị thứ nhất của mẫu số liệu ghép nhóm trên là bao nhiêu? \textit{(kết quả làm tròn đến hàng đơn vị)}.
    \shortans{164}
    \loigiai{
        Cỡ mẫu $N = 35$. Vị trí $Q_1$ là $\frac{N}{4} = 8,75$.
        Tần số tích lũy đến nhóm 1 là $6$, nhóm 2 là $18$. Nhóm chứa $Q_1$ là $[163;166)$.
        $Q_1 = 163 + \frac{8,75 - 6}{12} \cdot 3 = 163,6875$. Làm tròn đến hàng đơn vị ta được $164$.
    }
\end{ex}
\begin{ex}
    Thống kê về thời lượng mỗi trận đấu bi-a trong vòng tứ kết giải đấu European Open người ta được mẫu số liệu ghép nhóm như sau:
    \begin{center}
        \begin{tabular}{|c|c|c|c|c|c|}
            \hline
            Thời gian & $[9,5;12,5)$ & $[12,5;15,5)$ & $[15,5;18,5)$ & $[18,5;21,5)$ & $[21,5;24,5)$ \\
            \hline
            Số trận & $3$ & $12$ & $15$ & $24$ & $2$ \\
            \hline
        \end{tabular}
    \end{center}
    Số trung bình của mẫu số liệu trên là \textit{(kết quả làm tròn đến hàng phần mười)}.
    \shortans{17,5}
    \loigiai{
        Cỡ mẫu $N = 56$.
        Giá trị đại diện các nhóm lần lượt là: $11, 14, 17, 20, 23$.
        $\bar{x} = \frac{11 \cdot 3 + 14 \cdot 12 + 17 \cdot 15 + 20 \cdot 24 + 23 \cdot 2}{56} \approx 17,535$.
        Làm tròn đến hàng phần mười ta được $17,5$.
    }
\end{ex}
\begin{ex}
    Thống kê về thời lượng mỗi trận đấu bi-a trong vòng tứ kết giải đấu European Open người ta được mẫu số liệu ghép nhóm như sau:
    \begin{center}
        \begin{tabular}{|c|c|c|c|c|c|}
            \hline
            Thời gian & $[9,5;12,5)$ & $[12,5;15,5)$ & $[15,5;18,5)$ & $[18,5;21,5)$ & $[21,5;24,5)$ \\
            \hline
            Số trận & $3$ & $12$ & $15$ & $24$ & $2$ \\
            \hline
        \end{tabular}
    \end{center}
    Tứ phân vị thứ ba của mẫu số liệu trên là
    \shortans{20}
    \loigiai{
        Cỡ mẫu $N = 56$. Vị trí $Q_3$ là $\frac{3N}{4} = 42$.
        Tần số tích lũy đến nhóm 3 là $30$, đến nhóm 4 là $54$. Nhóm chứa $Q_3$ là $[18,5;21,5)$.
        $Q_3 = 18,5 + \frac{42 - 30}{24} \cdot 3 = 18,5 + 1,5 = 20$.
    }
\end{ex}
\begin{ex}
    Kết quả điều tra chiều cao của học sinh (đơn vị: cm) trong trường THPT được biểu diễn ở biểu đồ sau
    \begin{center}
        \begin{tikzpicture}[>=stealth, yscale=0.4,xscale=1]
            % Vẽ các trục tọa độ
            \draw[->] (0,0) -- (13,0) node[below] {(cm)};
            \draw[->] (0,0) -- (0,13.5) node[left] {(Số học sinh)};
            \node[below left] at (0,0) {$O$};
            % Ký hiệu ngắt trục (zigzag) ở gần gốc tọa độ
            \draw (0.4,-0.15) -- (0.5,0.15) -- (0.6,-0.15) -- (0.7,0.15) -- (0.8,0);
            \draw[dotted] (0.8,0) -- (1.5,0); 
            \foreach \x in {0.2, 0.4} \draw (\x,0.05) -- (\x,-0.05);
            % Vẽ các cột của biểu đồ (Mở rộng độ rộng mỗi cột lên 2 đơn vị)
            \draw[pattern=north east lines] (1.5,0) rectangle (3.5,6);
            \draw[pattern=dots] (3.5,0) rectangle (5.5,12);
            \draw[pattern=north west lines] (5.5,0) rectangle (7.5,9);
            \draw[pattern=horizontal lines] (7.5,0) rectangle (9.5,5);
            \draw[pattern=crosshatch] (9.5,0) rectangle (11.5,3);
            % Đường gióng nét đứt từ trục tung sang các cột
            \draw[dashed] (0,3) node[left] {$3$} -- (9.5,3);
            \draw[dashed] (0,5) node[left] {$5$} -- (7.5,5);
            \draw[dashed] (0,6) node[left] {$6$} -- (1.5,6);
            \draw[dashed] (0,9) node[left] {$9$} -- (5.5,9);
            \draw[dashed] (0,12) node[left] {$12$} -- (3.5,12);
            % Nhãn các nhóm trên trục hoành (Đã căn lại tọa độ tâm mỗi cột)
            \node[below] at (2.5,-0.1) {$[160;163)$};
            \node[below] at (4.5,-0.1) {$[163;166)$};
            \node[below] at (6.5,-0.1) {$[166;169)$};
            \node[below] at (8.5,-0.1) {$[169;172)$};
            \node[below] at (10.5,-0.1) {$[172;175)$};
        \end{tikzpicture}
    \end{center}
    Mốt của mẫu số liệu ghép nhóm trên là bao nhiêu?
    \shortans{165}
    \loigiai{
        Nhóm có tần số lớn nhất là $[163;166)$ với tần số $12$.
        $M_o = 163 + \frac{12 - 6}{(12 - 6) + (12 - 9)} \cdot 3 = 163 + \frac{6}{9} \cdot 3 = 165$.
    }
\end{ex}
\begin{ex}
    Kết quả khảo sát độ cận thị mắt của 100 học sinh trường tiểu học X được cho như bảng sau:
    \begin{center}
        \begin{tabular}{|c|c|c|c|c|c|c|}
            \hline
            Độ cận & $[0,25;0,75)$ & $[0,75;1,25)$ & $[1,25;1,75)$ & $[1,75;2,25)$ & $[2,25;2,75)$ & $[2,75;3,25)$ \\
            \hline
            Số học sinh & $25$ & $38$ & $21$ & $12$ & $3$ & $1$ \\
            \hline
        \end{tabular}
    \end{center}
    Bác sĩ sẽ chọn ra $25\%$ học sinh có độ cận cao nhất để tư vấn phụ huynh trong vấn đề chăm sóc đôi mắt, giảm khả năng tăng độ nhanh. Học sinh có độ cận từ bao nhiêu độ thì bác sĩ sẽ mời tư vấn phụ huynh? \textit{(kết quả làm tròn đến hàng phần trăm)}.
    \shortans{1,54}
    \loigiai{
        Để tìm mốc cho $25\%$ học sinh có độ cận cao nhất, ta cần tìm tứ phân vị thứ ba $Q_3$.
        Cỡ mẫu $N = 100$. Vị trí của $Q_3$ là $75$.
        Tần số tích lũy đến nhóm 2 là $25 + 38 = 63$, đến nhóm 3 là $63 + 21 = 84$. Nhóm chứa $Q_3$ là $[1,25;1,75)$.
        $Q_3 = 1,25 + \frac{75 - 63}{21} \cdot 0,5 \approx 1,5357$.
        Làm tròn đến hàng phần trăm, ta được $1,54$. Học sinh có độ cận từ $1,54$ độ trở lên sẽ được mời.
    }
\end{ex}
\Closesolutionfile{ans}